Due to the four-particle final states and the unavailability of $K/\pi$ separation,
there is a significant number of events with multiple candidates.

Hence we can include two further stages of selection in order to
obtain a single candidate per event, which can be useful for various tests.
As mentioned above, each event consists of multiple \Bd~candidates, and each candidate has
two possible mass hypotheses, i.e., whether an individual candidate corresponds to a \Bd or \Bdb.
To select amongst different four-particle sets, the single best four-particle set per event is
chosen based on the best $\chi^2$/Nd.o.f., referred to as the best-$\chi^2$ selection.

After the best-$\chi^2$ selection, we can still have two candidates corresponding to the two
hadron assigments ($K\pi$ and $\pi K$), Hence we can assign a single reconstructed mass to
the selected four-track set choosing between the two mass hypotheses (as the switching of
the $K$ and $\pi$ tracks lead to a change in the reconstructed $B$ mass).
The assignment of the $B$ mass is done in two ways: a comparison to the $K^*$ mass PDG value
is done and the closer of the two hypotheses is chosen, or in the gNN architecture,
a multi-class gNN produces a classification score for each candidate as \Bd, \Bdb, and background
and the highest of these is chosen.
The former of these is referred to as the best-$K^*$ mass selection.

The two additional selections above (the best-$\chi^2$ and the best-$K^*$ mass) are labelled
the best-candidate selection in the following and in the cutflow
tables~\ref{tab:presel:sigele},~\ref{tab:presel:resele},~\ref{tab:presel:sigmu} and~\ref{tab:presel:resmu}.


The percentage of events with more than one four-particle sets before the preselection is about $84.2\%$ ($82.1\%$)
for the non-resonant (resonant) sample for the electron final state,
while the same percentage after the preselection is down to about $18.7\%$ ($23.3\%$)
for the non-resonant (resonant) sample for the electron final state.

The percentage of events with more than one candidate before the preselection is about $73.2\%$ ($71.8\%$)
for the non-resonant (resonant) sample for the muon final state,
while the same percentage after the preselection is down to about $46.5\%$ ($42.1\%$)
for the non-resonant (resonant) sample for the muon final state.

\begin{table}[h]
\begin{center}
\begin{tabular}{cccccc}
\hline
            &  \% of events with           &  \% of events with  \\
            &  $\geq 1$ four-particle sets &  $\geq 1$ candidate \\
\hline\hline
electron final states   & & \\
\hline
~~~~non-resonant signal & $18.7\%$  & \\
~~~~resonant $J/\psi K^*$ & $23.3\%$ & \\
\hline\hline
muon final states       & & \\
\hline
~~~~non-resonant signal & & $46.5\%$ \\
~~~~resonant $J/\psi K^*$ &  & $42.1\%$ \\
\hline\hline
\end{tabular}
\end{center}
\caption{Percentages of multiple candidates after pre-selection in both the non-resonant and
resonant MC samples.}
\label{tab:multiplecandidates}
\end{table}









